%%
%% This is file `zjureport-example.tex',
%% generated with the docstrip utility.
%%
%% The original source files were:
%%
%% zjureport.dtx  (with options: `example')
%% This is a generated file.
%% Copyright (C) 2026 by Your Name
\documentclass[codeformat]{zjureport}
\usepackage{hyperref}
\hypersetup{colorlinks=true}

\major{信息工程}
\name{张扬}
\stuid{3171002333}
\college{信息与电子工程学院}
\date{\zhtoday}
\lab{东四-A-308}
\course{\LaTeX 模板使用}
\instructor{赵林平}
\grades{93}
\expname{使用实验报告模板}
\exptype{操作实验}
\partner{无}

\begin{document}
\makecover
\makeheader
\section{实验目的和要求}
测试楷体与非粗体封面内容效果。

\section{实验内容和步骤}

\subsection{实验内容}

\subsection{实验步骤}

\section{主要仪器设备}
计算机及配套软件。

\section{操作方法和实验步骤}

\section{实验数据记录和处理}

\begin{lstlisting}[language=C]
#include <stdio.h>
int main() {
    int number;

    printf("Enter an integer: ");

    // reads and stores input
    scanf("%d", &number);

    // displays output
    printf("You entered: %d", number);

    return 0;
}
\end{lstlisting}

\section{实验结果与分析}
\end{document}
\endinput
%%
%% End of file `zjureport-example.tex'.
